\worksheettitle
  {Класс миллиардов}
  {\modulemark{M03\(\star\)} Усложнение}

\studentnameblock

\begin{problembox}[Четыре класса]
  Рассмотрите число:
  \[
    7\,|\,304\,|\,052\,|\,018.
  \]

  \begin{center}
  \renewcommand{\arraystretch}{1.45}
  \begin{tabular}{c|cccc}
    \toprule
    класс & миллиарды & миллионы & тысячи & единицы \\
    \midrule
    группа & 7 & 304 & 052 & 018 \\
    число класса & \answerblank[14mm] & \answerblank[16mm] &
      \answerblank[16mm] & \answerblank[16mm] \\
    \bottomrule
  \end{tabular}
  \end{center}
\end{problembox}

\begin{practicebox}[Восстановите число]
  Класс миллиардов содержит \(4\), класс миллионов --- \(12\), класс тысяч
  --- \(5\), класс единиц --- \(8\).

  Запишите группы, сохраняя три позиции справа от первого класса:
  \[
    4\;|\;\answerblank[18mm]\;|\;\answerblank[18mm]\;|\;\answerblank[18mm].
  \]

  Всё число: \answerblank[55mm].
\end{practicebox}

\begin{warningbox}[Проверьте нулевые позиции]
  Числа классов \(12\), \(5\) и \(8\) записываются группами \(012\),
  \(005\) и \(008\), потому что каждый непервый класс занимает три разряда.
\end{warningbox}

\begin{reflectionbox}[Обобщите]
  Почему только самый левый класс может быть записан одной или двумя цифрами?
  \writinglines{2}
\end{reflectionbox}

